\documentstyle[%


righttag%


,11pt%


,amssymb%


,russian%               remove in Eng.


,izvru%                 Set izven% in Eng.


]{amsart}





% Move the footnote to \thanks in Eng.


% {\in}


\newcommand{\const}{\operatorname{const}}


\newcommand{\mes}{\operatorname{mes}}


\newcommand{\tg}{\operatorname{tg}}


\newcommand{\gea}{\geqslant}


\newcommand{\lea}{\leqslant}


\newcommand{\vp}{\varphi}


\newcommand{\si}{\sigma}


\newcommand{\vk}{\varkappa}


\newcommand{\ve}{\varepsilon}


\newcommand{\N}{\Bbb N}


\newcommand{\R}{\Bbb R}


\newcommand{\rL}{\cal L}


\newcommand{\intl}{\int\limits}


\newcommand{\suml}{\sum\limits}


\newcommand{\ol}{\overline}


\newcommand{\sgn}{\operatorname{sgn}}


\newcommand{\ctg}{\operatorname{ctg}}


\newcommand{\lbr}{[\hskip -0.5mm [} %% pay attn. on final form


\newcommand{\rbr}{]\hskip -0.5mm ]}


\numberwithin{equation}{section}


\newcommand{\tts}[1]{}








%\hoffset=-15mm%                  For Eng.


\setcounter{page}{11}


\god{2004}


\nomer{8~(507)} %                        Remove in Eng.


%\nomer{Vol.\,48, No.\,8}%               Set in Eng.


%\str{\pageref{begin}--\pageref{end}}%  Set in Eng.


\udk{517.968:519.642}





\author{\it Б.Г.\,Габдулхаев}


% NB:   ^^ remove in Eng.


\title{Методы решения бисингулярных интегральных уравнений


с внутренними коэффициентами}





\date{}


%\pagestyle{myheadings}%         Set in Eng.


\pagestyle{plain} %              Remove in Eng.


\begin{document}


%\thispagestyle{plain}%          Set in Eng.


\maketitle


\label{begin}











\section*{Введение}


\setcounter{section}{0}





Значительное число теоретических и практических задач приводит к необходимости 


решения следующего двумерного сингулярного интегрального уравнения (с.\,и.\,у.) 


с внутренними коэффициентами:


\begin{multline}


A\vp \equiv a_0(s,\si)\vp(s,\si)+\frac{1}{2\pi}\int_{0}^{2\pi}a_1(s,\si;\xi)


\ctg\frac{\xi-s}{2}\vp(\xi,\si)d\xi + \\


%


+\frac{1}{2\pi}\int_{0}^{2\pi}a_2(s,\sigma;\eta)


\ctg \frac{\eta-\sigma}{2}\vp(s,\eta)d\eta + \\


%


+\frac{1}{4\pi^2}\int_0^{2\pi}\int_0^{2\pi}


a_{12}(s,\sigma;\xi,\eta)\ctg\frac{\xi-s}{2}


\ctg\frac{\eta-\sigma}{2}\vp(\xi,\eta)d\xi\,d\eta=


f(s,\si),\ \ -\infty<s,\sigma<\infty.


\end{multline}


Здесь $a_0(s,\si)$, $a_1(s,\si;\xi)$, $a_2(s,\si;\eta)$, 


$a_{12}(s,\si;\xi,\eta)$ --- известные непрерывные $2\pi$-периодические


функции\footnote{Они таковы, что оператор $A$ из (0.1) ограничен в 


$L_2[0,2\pi]^2$. Для этого достаточно, чтобы функции $a_i$ ($i=1,2,12$) были 


непрерывны по Г\"eльдеру; если же $a_1(s,\si;\xi)=b_1(s,\si)c_1(\xi)$, 


$a_2(s,\sigma;\eta)=b_2(s,\si)c_2(\eta)$ и 


$a_{12}(s,\si;\xi,\eta)=b_{12}(s,\si)c_{12}(\xi,\eta)$, то достаточно, 


чтобы функции $b_i$ и $c_i$ ($i=1,2,12$) были непрерывными в своих областях 


определения.} по каждой из переменных, 


$f(s,\si)$ и $\vp(s,\si)\in L_2[0,2\pi]^2\equiv L_2$ ---


данная и искомая функции соответственно, причем сингулярные интегралы 


\begin{align}


S_1(a_1\varphi)&=S_1(a_1\varphi ;s,\sigma )=\frac{1}{2\pi}\int_0^{2\pi}


a_1(s,\sigma ;\xi )\ctg\frac{\xi -s}{2} \varphi (\xi ,\sigma )d\xi ,\\


%


S_2(a_2\varphi)&=S_2(a_2\varphi ;s,\sigma )=\frac{1}{2\pi}\int_0^{2\pi}


a_2(s,\sigma ;\eta )\ctg \frac{\eta -\sigma }{2}


\varphi (s,\eta )d\eta ,\\


%


S_{12}(a_{12}\varphi)&=


S_{12}(a_{12}\varphi ;s,\sigma )


=\frac{1}{4\pi^2}


\int_0^{2\pi}\int_0^{2\pi}


a_{12}(s,\sigma ;\xi ,\eta )


\ctg \frac{\xi -s}{2}\ctg \frac{\eta -\sigma }{2}


\varphi (\xi ,\eta ) d\xi\, d\eta


\end{align}


понимаются в смысле главного значения по Коши--Лебегу [1]--[3].





Теория уравнений вида (0.1), в том числе аппроксимативных методов их решения,


весьма сложна и все еще далека от своего завершения.


Ниже, в продолжение ряда результатов


[4]--[8], приводятся простые достаточные условия существования,


единственности и устойчивости решений уравнения (0.1), а также предлагаются


практически эффективные (в том числе оптимальные) аппроксимативные методы


его решения. Ради простоты


выкладок все функции и пространства считаем (без ограничения общности)


вещественными. Отметим также, что некоторые результаты работы анонсированы 


в ([9]; [10], гл.\,4, \S\,4).








\section{Основные результаты}





1.1. {\it Теорема существования и единственности решения.} В следующей 


теореме устанавливается существование, единственность  и устойчивость 


решения уравнения (0.1).





\begin{thm}


Пусть выполнены условия


$|a_0(s,\sigma )| \gea \alpha=\const > 0$,


$a_1(s,\sigma ;\xi )=a_1(\xi ,\sigma ;s)$, 


$a_2(s,\sigma ;\eta )=a_2(s,\eta ;\sigma)$,


$a_{12}(s,\sigma ;\xi ,\eta )=-a_{12}(\xi ,\eta ;s,\sigma )$. 





Тогда оператор $A: L_2\to L_2$ непрерывно обратим и 


$$


\|A^{-1}\|\lea \alpha^{-1}<\infty,


$$


а уравнение $(0.1)$ имеет единственное решение 


$\varphi^*(s,\sigma ){\in} L_2$ при любой правой части


$f(s,\sigma ){\in} L_2$, которое устойчиво относительно малых


возмущений $f(s,\sigma )$ в пространстве $L_2$ $($и тем более


в пространстве $C[0,2\pi]^2)$.


\end{thm}





\begin{pf}


Обозначим через $L_2\equiv L_2[0,2\pi]^2$ пространство квадратично 


суммируемых в $[0,2\pi;0,2\pi]\equiv [0,2\pi]^2$ по Лебегу функций c 


обычными скалярным произведением и нормой соответственно


\begin{alignat*}{2}


(f,g)&=\frac{1}{4\pi^2}\int_{0}^{2\pi}\int_{0}^{2\pi}


f(s,\si)g(s,\si)ds\, d\si &\quad (f,g&\in L_2),\\


%


\|f\|&=\bigg\{


\frac{1}{4\pi^2}\int_{0}^{2\pi}\int_{0}^{2\pi}|f(s,\si)|^2 ds\, d\si


\bigg\}^{1/2} &\quad (f&\in L_2).


\end{alignat*}


Через $C = C[0,2\pi]^2$ будем обозначать пространство непрерывных 


$2\pi$-периодических по каждой из переменных функций от двух переменных 


с обычной нормой 


$$


\|g\|_{C[0,2\pi]^2} = \max\limits_{-\infty<s,\sigma<\infty}|g(s,\sigma)|


\quad (g\in C[0,2\pi]^2).


$$


Тогда с.\,и.\,у. (0.1) эквивалентно заданному в $L_2$ линейному 


операторному уравнению 


\begin{equation}


A\vp \equiv S_0(a_0\vp) + S_1(a_1\vp)+S_2(a_2\vp)+ S_{12}(a_{12}\vp) = f


\quad (\vp, f\in L_2),


\end{equation}


где $S_0(a_0\vp)=a_0(s,\si)\vp(s,\si)$, причем


$A \equiv S_0a_0 + S_1 a_1 + S_2 a_2 + S_{12}a_{12}:L_2\to L_2$ 


есть непрерывный оператор:


$$


\|A\|_{L_2\to L_2}\lea \suml_{i=0,1,2,12}\|S_ia_i\|_{L_2\to L_2}\lea M


= \const < \infty.


$$





Далее без ограничения общности будем считать, что $a_0(s,\sigma)>0$; 


случай $a_0(s,\sigma)<0$ рассматривается аналогично.





Очевидно, для любой функции $\vp \in L_2$ имеем


\begin{equation}


(S_0(a_0\vp),\vp)=\frac{1}{4\pi^2}\int_{0}^{2\pi}\int_{0}^{2\pi}


a_0(s,\si)|\vp(s,\si)|^2 ds\, d\si \gea \alpha\|\vp\|^2.


\end{equation}


В силу условий на функции $a_i$ ($i=1,2,12$) и свойств ядер Гильберта 


для любой функции $\vp \in L_2$ находим


\begin{equation}


K_i\equiv (S_i(a_i\vp),\vp)=0\quad (i=1,2,12).


\end{equation}


Действительно, для интеграла (0.2) имеем 


\begin{gather*}


K_1\equiv (S_1(a_1\vp),\vp)= \frac{1}{4\pi^2}\int_{0}^{2\pi}\int_{0}^{2\pi}


\vp(s,\si)ds\,d\si \frac{1}{2\pi} \int_{0}^{2\pi}


a_1(s,\si;\xi)\ctg\frac{\xi-s}{2}\vp(\xi,\si)d\xi = \\


%


%= \frac{1}{8\pi^3}\int_{0}^{2\pi}d\sigma \int_{0}^{2\pi}\vp(s,\sigma) ds


%\int_{0}^{2\pi} a_1(s,\sigma;\xi)\ctg\frac{\xi-s}{2}\vp(\xi,\si)d\xi=\\


%


=\frac{1}{8\pi^3}\int_0^{2\pi}d\si\int_{0}^{2\pi}\vp(\xi,\si)d\xi


\int_0^{2\pi}a_1(\xi,\si;s)\ctg\frac{s-\xi}{2}\vp(s,\si)ds= \\


%


=\frac{1}{8\pi^3}\int_{0}^{2\pi}d\si \int_0^{2\pi}\vp(s,\si)ds


\int_0^{2\pi}a_1(s,\si;\xi)\bigg[-\ctg\frac{\xi-s}{2}\bigg]


\vp(\xi,\sigma)d\xi = -K_1,


\end{gather*}


поэтому $K_1=0$. Аналогично для интеграла (0.3) получаем 


$K_2=(S_2(a_2\vp),\vp)=0$ для любой $\vp \in L_2$. С учетом свойств функций 


$a_{12}$ и ядер Гильберта для интеграла (0.4) при любых $\vp \in L_2$ находим


\begin{gather*}


K_{12} \equiv (S_{12}(a_{12}\vp),\vp) = \\


%


=\frac{1}{4\pi^2}\int_{0}^{2\pi}\int_0^{2\pi}\vp(s,\si)ds\,d\si


\frac{1}{4\pi^2}\int_0^{2\pi}\int_0^{2\pi} a_{12}(s,\si;\xi,\eta)


\ctg\frac{\xi-s}{2}\ctg\frac{\eta-\si}{2}\vp(\xi,\eta)d\xi\,d\eta= \\


%


=\frac{1}{16\pi^4}\int_0^{2\pi}\int_{0}^{2\pi} \vp (\xi,\eta)d\xi\,d\eta


\int_0^{2\pi}\int_0^{2\pi} a_{12}(\xi,\eta;s,\si)\ctg\frac{s-\xi}{2}


\ctg\frac{\si-\eta}{2}\vp(s,\si)ds\, d\si = \\


%


= \frac{1}{16\pi^4} \int_0^{2\pi} \int_0^{2\pi} \vp(s,\si) ds\,d\si


\int_0^{2\pi}\int_0^{2\pi} [-a_{12}(s,\si;\xi,\eta)]\ctg\frac{\xi-s}{2}


\ctg\frac{\eta-\si}{2}\vp(\xi,\eta)d\xi\,d\eta = -K_{12},


\end{gather*}


следовательно, $K_{12} = 0$.





Из соотношений (1.1)--(1.3) находим неравенство


\begin{equation}


(A\vp,\vp) \gea \alpha\|\vp\|^{2},\quad \vp \in L_2.


\end{equation}


Аналогично для оператора $A^*:L_2\to L_2$, сопряженного к оператору 


$A:L_2\to L_2$ в пространстве $L_2$, имеем


\begin{equation}


(A^*\vp,\vp) = (\vp,A\vp) = (A\vp,\vp) \gea \alpha\|\vp\|^2,\quad \vp\in L_2.


\end{equation}





Из соотношений (1.4) и (1.5) находим соответственно неравенства 


$$


\|A\vp\|\gea \alpha\|\vp\|,\quad


\|A^*\vp\|\gea \alpha\|\vp\|,\quad \vp \in L_2,


$$


обеспечивающие [11] существование левых ограниченных обратных операторов 


$A_l^{-1}$ и $(A^*)_l^{-1}$ соответственно


$$


\|A^{-1}_l\|\lea \alpha^{-1}<\infty \quad \text{и}


\quad \|(A^*)^{-1}_l\|\lea \alpha^{-1}<\infty.


$$


Отсюда, в свою очередь, следует [11] существование и ограниченность 


двустороннего обратного оператора $A^{-1}:L_2\to L_2$ и справедливость 


необходимых для него неравенств.





Остальные утверждения теоремы легко выводятся из сказанного выше.


\end{pf}





1.2. {\it Итерационный метод.} В следующей теореме устанавливается 


возможность решения уравнения (0.1) итерационным методом.





\begin{thm}


В условиях теоремы $1$ решение $\varphi^* \in L_2$ 


уравнения $(0.1)$ при любой правой части $f\in L_2$ может быть найдено 


с помощью линейного


универсального итерационного метода


\begin{equation}


\varphi ^k=\varphi ^{k-1}+\alpha \|A\|^{-2}(f-A\varphi ^{k-1}),\quad


k=1,2,\dotsc ,


\end{equation}


при любом начальном приближении $\varphi^0\in L_2$. Погрешность


$k$-го приближения может быть оценена в пространстве


$L_2$ неравенствами


\begin{equation}


\|\varphi ^\ast -\varphi ^k\|\lea q^k\|\varphi ^\ast -\varphi ^0\|\lea


q^k(1-q)^{-1}\|\varphi ^1-\varphi ^0\| \quad (k=1,2,\dotsc),


\end{equation}


где $q=(1-\alpha ^2\|A\|^{-2})^{1/2}<1;$ если же начальное приближение 


выбирается по формуле $\varphi ^0=


\alpha \|A\|^{-2}f$, то и неравенствами


\begin{equation}


\|\varphi ^\ast -\varphi ^k\|\lea q^{k+1}(1-q)^{-1}


\alpha \|A\|^{-2}\|f\|,\quad k\in \N.


\end{equation}


\end{thm} 





\begin{pf}


Уравнение (0.1) эквивалентно заданному в $L_2$  линейному операторному 


уравнению второго рода вида


\begin{equation}


\vp = (E-\tau A)\vp + \tau f,\quad


\tau = \frac{\alpha}{\|A\|^2},\quad \vp \in L_2,


\end{equation}


где оператор перехода $B = E-\tau A$ является сжимающим оператором с 


коэффициентом сжатия $q = (1-\alpha/\|A\|^2)^{1/2}<1$. С другой стороны, 


универсальный итерационный метод (1.6) для уравнения (0.1) является 


методом простой итерации для эквивалентного ему уравнения (1.9). Поэтому 


утверждения теоремы, в том числе оценки (1.7) и (1.8), легко выводятся 


из известных результатов (напр., [11], с.\,213--214; [12], сс.\,104, 131) 


по итерационным методам решения операторных уравнений второго рода в 


полных линейных нормированных пространствах. 


\end{pf}





1.3. {\it Общий проекционный метод и его частные случаи.} Обозначим 


через $\{\psi_r=\psi_r(s,\sigma )\}_1^\infty$


полную ортонормальную систему функций в $L_2$, а через $c_r(f)=(f,\psi_r)$ ---


коэффициенты Фурье функции $f\in L_2$ по указанной системе.


Приближения $\varphi_N =\varphi_N (s,\sigma )$ к решению


$\varphi^\ast =\varphi^\ast (s,\sigma )$ уравнения (0.1) будем


искать в виде обобщенного ``полинома''


\begin{equation}


\varphi_N (s,\sigma )=\sum_{k=1}^N \alpha _k\, \psi_k(s,\sigma ),


\quad N=1,2,\dotsc,


\end{equation}


неизвестные коэффициенты $\alpha _k=\alpha _{k,N}$ которого будем


определять из системы линейных алгебраических уравнений (СЛАУ)


\begin{equation}


\sum_{k=1}^N\alpha _k\, c_r(A\psi_k)=c_r(f),\quad r=\overline{1,N}.


\end{equation}








\begin{thm} 


В условиях теоремы $1$ СЛАУ $(1.11)$ однозначно разрешима при любых


$N=1,2,\dotsc$, а приближенные решения $(1.10)$ сходятся к решению уравнения


$(0.1)$ в пространстве $L_2$, причем для погрешности справедливы


двусторонние оценки


\begin{equation}


E_N (\varphi^*)\lea \|\varphi^* -\varphi _N\|\lea


\alpha ^{-1}\|A\|E_N (\varphi^*),\quad N=1,2,\dotsc,


\end{equation}


где $E_N(g)$ --- наилучшее приближение функции $g\in L_2$


всевозможными элементами вида $(1.10)$ в пространстве $L_2=L_2[0,2\pi]^2$.


\end{thm}





\begin{pf}


Пусть $X = L_2 = L_2[0,2\pi]^2$, а $X_N = \{\vp_N\}$ --- 


его $N$-мерное подпространство, состоящее из всевозможных элементов 


вида (1.10). Тогда СЛАУ (1.11) эквивалентна операторному уравнению 


\begin{equation}


A_N\vp_N \equiv P_N A \vp_N = P_N f\quad  (\vp_N, P_N f\in X_N),


\end{equation}


где оператор проектирования $P_N:X\to X_N$ определяется по формуле


\begin{equation}


P_N(g;s,\si) = \sum_{r=1}^N c_r(g)\psi_r(s,\si),\quad g\in L_2.


\end{equation}


Нетрудно видеть, что 


\begin{equation}


P_N^2 = P_N,\quad P_N^* = P_N,\quad \|P_N\|_{X\to X_N \subset X}=1,


\quad N\in \N.


\end{equation}


Поэтому в силу (1.4), (1.13) и (1.14) для любых $\vp_N\in X_N$ находим


\begin{equation}


(A_N\vp_N,\vp_N) = (P_NA\vp_N,\vp_N) = (A\vp_N,P_N^*\vp_N)=


(A\vp_N,P_N\vp_N) = (A\vp_N,\vp_N) \gea \alpha\|\vp_N\|^2.


\end{equation}


Отсюда следует неравенство


\begin{equation}


\|A_N\vp_N\|\gea \alpha\|\vp_N\|,\quad \vp_N \in X_N,


\end{equation}


обеспечивающее (напр., [13], гл.\,I, \S\,2) двустороннюю обратимость 


оператора $A_N:X_N\to X_N$ при любых $N\in \N$ и справедливость оценки 


\begin{equation}


\|A_N^{-1}\|\lea \alpha^{-1}<\infty,\quad N\in\N.


\end{equation}


Поэтому как уравнение (1.13), так и эквивалентная ему СЛАУ (1.11) 


однозначно разрешимы при любых натуральных $N$ и любых правых частях, 


причем приближенное решение (1.10) в силу (1.15)--(1.18) удовлетворяет


неравенствам


$$


\|\vp_N\| = \|A_N^{-1}P_N f\|\lea \alpha^{-1}\|P_N f\|\lea \alpha^{-1}\|f\|,


\quad N\in \N.


$$





Для решений $\vp^* = A^{-1}f$ и $\vp _N = A^{-1}_NP_N f$ уравнений 


соответственно (1.1) и (1.13) из теоремы~6 ([13], гл.\,I) следует тождество


\begin{equation}


\vp^* - \vp_N = (E-A_N^{-1}P_N A)(\vp^* - \ol{\vp}_N),


\end{equation}


где $E$ --- единичный оператор, а $\ol{\vp}_N$ --- произвольный элемент 


из $X_N$. Полагая $\ol{\vp}_N = P_N \vp^*,$ из (1.15) и (1.19) находим 


неравенство 


\begin{equation}


\|\vp^* - \vp_N\|_X \lea \|E - A_N^{-1}P_N A\|_{X\to X} E_N(\vp ^*),


\quad N\in \N.


\end{equation}


Поскольку для любых $N\in \N$


\begin{equation}


(E - A_N^{-1}P_NA)^2 = E-A_N^{-1}P_NA,


\end{equation}


то с помощью соотношений (1.15) и (1.18) можно доказать


\begin{equation}


\|E - A_N^{-1}P_NA\| = \|A_N^{-1}P_NA\| \lea \alpha^{-1}\|A\|,\quad


N \in\N.


\end{equation}


Из (1.20) и (1.22) следуют оценки (1.12). Так как в условиях теоремы 


$E_N(g)\to 0$, $N\to \infty$, для любой функции $g\in X$, то сходимость 


метода доказана.


\end{pf}





На практике большое значение имеет выбор координатной системы функций


$\{ \psi_r(s,\sigma )\}$. Дадим несколько способов


такого выбора, которые приводят к конкретным проекционным методам 


решения с.\,и.\,у. (0.1). 


\medskip





1.3.1. {\it Метод редукции $($Гал\"еркина$)$.} 


Пусть $N=(2n+1)(2m+1)$, где $n$ и $m=0,1,\dotsc$, а


$\psi_r=\psi_r(s,\sigma )=


e^{i(ks+j\sigma )}\equiv \psi_{kj}(s,\si)$, $r=\overline{1,N}$,


$k=\overline{-n,n}$, $j=\overline{-m,m}$, $i^2=-1$. Тогда


$$


\varphi_N (s,\sigma )=\sum_{r=1}^N\alpha _r\, \psi_r(s,\sigma )=


\sum_{k=-n}^n\sum_{j=-m}^m\beta_{kj}\, e^{i(ks+j\sigma )} ,


$$


где $\alpha_r$, $\beta_{kj}$ --- неизвестные коэффициенты.





Очевидно, что в этом случае имеем дело с 


{\it методом редукции $($Гал\"еркина$)$}


решения с.\,и.\,у. (0.1), СЛАУ которого имеет вид 


$$


\sum_{k=-n}^n\ \sum_{j=-m}^m \beta_{kj} c_{rl}(A\psi_{kj}) =


c_{rl}(f),\quad r=\ol{-n,n},\ \ l=\ol{-m,m},


$$


где


$$


c_{rl}(g) = \frac{1}{4\pi^2}\int_{0}^{2\pi}\int_{0}^{2\pi}g(s,\si)


e^{-i(rs+l\si)}ds\, d\si,\quad g\in L_2.


$$





1.3.2. {\it Метод сплайн-подобластей. } Рассмотрим случай 


$\psi_r=\psi_{r,N}(s,\sigma )$, $r=\overline{1,N}$,


$N=1,2,\dotsc$, где последовательность подпространств $X_N =


\rL(\{\psi_{k,N}\}_1^N)$, натянутых на элементы


$\psi_1,\dotsc,\psi_N$, предельно плотна в пространстве $L_2$.


За $\psi_r=\psi_{r,N}(s,\sigma)$ возьмем характеристические функции


областей $\Delta_r=\Delta_{r,N}\subset [0,2\pi]^2$, $r=\overline{1,N}$,


$N=1,2,\dotsc$, где плоская мера $\mes (\Delta_k\cap


\Delta_j)=0$ для любых $k$, $j=\overline{1,N}$, $k\ne j$,


$\mathop{\cup}\limits_{r=1}^N


\Delta_r=[0,2\pi ]^2$, причем


$$


\lim_{N\rightarrow \infty }\; \max_{1\leq j\leq N}\; \mes \Delta


_{j,N}=0.


$$


В частном случае $N=nm$ положим $\psi_{r,N}(s,\sigma )=


\psi_{k,n}(s)\psi_{j,m}(\sigma )$, $r=\overline{1,N}$, $k=\overline{1,n}$,


$j=\overline{1,m}$, где $\psi_{k,n}(s)$ и $\psi_{j,m}(\sigma )$ ---


характеристические функции интервалов соответственно $(s_{k-1},s_k]$,


$k=\overline{1,n}$, и $( \sigma _{i-1},\sigma _i]$, 


$i=\overline{1,m}$, а


\begin{equation}


s_j=\frac{2j\pi}{n},\ \ j=\ol{0,n};\quad \sigma _l=\frac{2l\pi}{m},


\ \ l=\ol{0,m}\quad


(n,m=1,2,\dotsc).


\end{equation}


Очевидно, что в этом случае имеем дело с {\it методом сплайн-подобластей 


нулевого порядка} решения


бисингулярного интегрального уравнения (0.1).





Заметим также, что в обоих частных случаях сходимость и оценки погрешности


указанных конкретных методов следуют из общей теоремы 3.





Теорема 3, в том числе оценки (1.12), показывают, что рассматриваемый в 


ней метод является {\it оптимальным по порядку} ([13], гл.\,II) среди 


всевозможных прямых и проекционных методов, позволяющих построить 


приближенное решение с.\,и.\,у. (0.1) в виде обобщенного ``полинома''


(1.10). Если же уравнения (1.1) и (1.13) таковы, что 


\begin{equation}


\|E - A_N^{-1}P_N A\|_{X\to X} \rightarrow 1,\quad N\to \infty,


\end{equation}


то исследуемый метод будет {\it асимптотически оптимальным} ([13], гл.\,II);


если же


\begin{equation}


\|E - A_N^{-1}P_N A\|_{X\to X} = 1,\quad N\in \N,


\end{equation}


то указанный метод будет {\it оптимальным} ([13], гл.\,II). Однако следует 


отметить, что условие (1.25) выполняется лишь при очень жестких 


ограничениях на коэффициенты уравнения (0.1), т.\,к. в силу соотношения 


(1.21) в общем случае имеем $\|E-A_N^{-1}P_N A\|\gea 1$ для любых $N\in\N$. 


Тем не менее класс уравнений (0.1), удовлетворяющих условию (1.25), не 


пуст\,; например, это условие выполняется в любом из следующих случаев:


\begin{itemize}


\item[1)] $a_k = \const$ ($k=0,1,2,12$), $a_0\ne 0$.


\item[2)] $a_0 = \const \ne 0$, $a_1 = a_1(s-\xi)$, $a_2 = a_2(\si-\eta)$,


$a_{12} = a_{12}(s-\xi,\si-\eta)$, 


где $a_1(t), a_2(\tau), a_{12}(t,\tau)$ --- непрерывные 


$2\pi$-периодические функции своих аргументов. 


\end{itemize}





Ясно, что  множество уравнений вида (0.1), удовлетворяющих условию (1.24), 


тем более не пусто. 





Поясним сказанное на примере метода редукции из п.\,1.3.1. 


Сначала рассмотрим случай~1). Поскольку здесь 


$$


A_N \vp_N \equiv P_NA\vp_N = A\vp_N = P_N f \quad (\vp_N, P_Nf\in X_N),


$$


то справедливы тождества


$$


A_N^{-1}P_NA(\vp^* - P_N\vp^*) = (A^{-1}P_N - P_NA^{-1})f =


(A^{-1}P_NA - P_N)A^{-1}f = A^{-1}(P_N - AP_NA^{-1})f,


$$


где $\vp^*$ и $\vp_N$ --- соответственно точное и приближенное решения 


с.\,и.\,у. (0.1). 





Для любой функции $g\in L_2$  имеем


$$


g(s,\si) = \sum_{k=-\infty}^\infty\ \sum_{r=-\infty}^\infty


c_{kr}(g)e^{i(ks+r\si)},\quad 


P_N(g;s,\si)=\sum_{k=-n}^n\ \sum_{r=-m}^m c_{kr}(g)e^{i(ks+r\si)},


$$


где $N = (2n+1)(2m+1)$.


Поэтому с.\,и.\,у. (0.1) эквивалентно уравнению 


$$


\sum_{k=-\infty}^\infty\ \sum_{r=-\infty}^\infty c_{kr}(\vp)


G_{kr}e^{i(ks+r\si)}=


\sum_{k=-\infty}^\infty\ \sum_{r=-\infty}^\infty c_{kr}(f)e^{i(ks+r\si)},


$$


где 


$G_{kr} = a_0+ia_1\sgn k + i a_2 \sgn r - a_{12}\sgn k \sgn r$,


$\sgn\beta = \{+1$ при $\beta>0$; $0$ при $\beta=0$; $-1$ при $\beta<0\}$.


Отсюда находим


$c_{kr}(\vp)=c_{kr}(f)/G_{kr}$ ($k=\ol{-\infty,\infty}$, $r=\ol{-\infty,\infty}$).


Поэтому решение с.\,и.\,у. (0.1) при любой $f\in L_2$ представимо в виде


$$


\vp^*(s,\si) = A^{-1}(f;s,\si) = \sum_{k=-\infty}^\infty\ \sum_{r=-\infty}^\infty


\frac{c_{kr}(f)}{G_{kr}}e^{i(ks+r\si)}.


$$


Отсюда находим


\begin{gather*}


P_N(A^{-1}f;s,\si) = \sum_{k=-n}^n\ \sum_{r=-m}^m \frac{c_{kr}(f)}{G_{kr}}


e^{i(ks+r\si)}, \\


%


A^{-1}(P_Nf;s,\si) = \sum_{k=-\infty}^\infty\ \sum_{r=-\infty}^\infty


\frac{c_{kr}(P_Nf)}{G_{kr}}e^{i(ks+r\si)}=


\sum_{k=-n}^n\ \sum_{r=-m}^m


\frac{c_{kr}(f)}{G_{kr}}e^{i(ks+r\si)} = P_N(A^{-1}f;s,\si);


\end{gather*}


при выводе последнего равенства учтено также, что 


$$


c_{kr}(P_Nf) =


\begin{cases}


c_{kr}(f) & \text{при } |k|\lea n,\ |r|\lea m;\\


0 &  \text{при } |k| \lea n,\  |r| > m;\\


0 &  \text{при } |k| > n,\  |r| \lea m;\\


0 &  \text{при } |k| > n,\ |r| > m,


\end{cases}


$$


где $k=\ol{-\infty,\infty}$, $r=\ol{-\infty,\infty}$, $N=(2n+1)(2m+1)$. 





Из приведенных соотношений и формулы (1.19) при $\ol{\vp}_N=P_N\vp^*$ 


следует, что соотношение (1.25) выполняется при любых $N\in\N$, 


следовательно, в рассматриваемом случае метод редукции является оптимальным. 





Такое же утверждение справедливо и в случае 2). Здесь существенно 


пользуемся соотношениями 


\begin{align*}


S_1(a_1\vp;s,\si) &= \sum_{k=-\infty}^\infty\ \sum_{r=-\infty}^\infty


c_{kr}(\vp) c_{k\infty}(a_1)e^{i(ks+r\si)},\\


%


S_2(a_2\vp;s,\si) &= \sum_{k=-\infty}^\infty\ \sum_{r=-\infty}^\infty


c_{kr}(\vp)c_{\infty r}(a_2)e^{i(ks+r\si)},\\


%


S_{12}(a_{12}\vp;s,\si) &= \sum_{k=-\infty}^\infty\ \sum_{r=-\infty}^\infty


c_{kr}(\vp)c_{kr}(a_{12})e^{i(ks+r\si)},


\end{align*}


где $\vp(s,\si)\in L_2$, $i=\sqrt{-1}$, а 


\begin{gather*}


c_{k\infty}(a_1) = \frac{1}{2\pi}\int_{0}^{2\pi}a_1(s)e^{-iks}ds,\quad 


c_{\infty r}(a_2) = \frac{1}{2\pi}\int_{0}^{2\pi}a_2(\si)e^{-ir\si}d\si,\\


%


c_{kr}(a_{12}) = \frac{1}{4\pi^2}\int_{0}^{2\pi}\int_0^{2\pi}


a_{12}(s,\si)e^{-i(ks+r\si)}ds\,d\si.


\end{gather*}


Поэтому собственными значениями оператора $A$ и соответствующими им 


собственными функциями будут 


$$


G_{kr} = a_0 + c_{k\infty}(a_1) + c_{\infty r}(a_2) + c_{kr}(a_{12}),\quad 


\psi_{kr}(s,\si)=e^{i(ks+r\si)},


$$


где $k=\ol{-\infty,\infty}$, $r = \ol{-\infty,\infty}$.





Доказательство завершается аналогично доказательству в случае 1).


\medskip





1.4. {\it Проекционно-итеративные методы.} Отметим, что даже 


при достаточно удачном выборе координатных функций


$\{\psi_r(s,\sigma )\}$ решение СЛАУ (1.11) во многих случаях представляет


значительные практические трудности. В связи с этим может оказаться интересным


следующий простой способ ее решения универсальным линейным одношаговым


итерационным методом:


\begin{equation}


\alpha _l^k=\alpha _l^{k-1}+\alpha \|A\|^{-2}\bigg\{ c_l(f)-


\sum_{r=1}^N \alpha_r^{k-1}c_l(A\psi_r)\bigg\},


\quad l=\overline{1,N};\ \ N,k=1,2,\dotsc,


\end{equation}


где $\alpha _1^0,\dotsc,\alpha_N^0$ --- произвольные вещественные числа.


Дело в том, что при $k\rightarrow \infty$ числовая последовательность


$\alpha _l^k$ ($l=\overline{1,N}$) сходится к решению


$\alpha _l$ ($l=\overline{1,N}$) СЛАУ (1.11) со скоростью


геометрической прогрессии с указанным выше знаменателем $q$, причем


для погрешностей\linebreak $|\alpha_l-\alpha _l^k|$, $l=\overline{1,N}$,


$k=1,2,\dotsc$, справедливы оценки типа (1.7) и (1.8).





Из теорем 1--3, как и в одномерном случае [14] 


(см. также [13], гл.\,II, \S\,7; [10], гл.\,4, \S\,1), выводится 





\begin{thm}


Проекционно-итеративная последовательность элементов


$$


\varphi_N^k(s,\sigma )=\sum_{r=1}^N\alpha _r^k\, \psi_r(s,\sigma ),\quad


k=1,2,\dotsc;\ \ N=1,2,\dotsc,


$$


сходится в $L_2$ к точному и приближенному решениям соответственно


$\, \varphi ^\ast (s,\sigma )\, $ и $\, \varphi_N (s,\sigma )\, $


уравнения $(0.1)$ в том смысле, что


$$


\varphi ^\ast =\lim_{N\rightarrow \infty }\, \varphi_N =


\lim_{N\rightarrow \infty}\; \lim_{k\rightarrow \infty}\, \varphi_N ^k,


$$


причем существует такой номер итерации $k=k_0(N)$, что для любых


$k\geq k_0(N)$


$$


\lim_{N\rightarrow\infty }\, \varphi_N ^k=\varphi ^\ast


$$


со скоростью геометрической прогрессии со знаменателем $\, q\, $,


указанным в теореме~$2$.


\end{thm}





1.5. {\it Квадратурно-кубатурные методы.} Результаты, аналогичные 


теоремам 3 и 4, справедливы также


для разработанного автором (см., напр.,  [4]--[8] и литературу в них)


квадратурно--кубатурного метода ре\-ше\-ния уравнения (0.1) на сетке узлов 


(1.23). Здесь приведем лишь некоторые результаты, основанные на 


аппроксимации тригонометрическими многочленами и полиномиальными 


сплайнами минимальных степеней.  


\medskip





1.5.1. {\it Полиномиальные квадартурно-кубатурные методы.} 


Приближенное решение с.\,и.\,у. (0.1) с непрерывными коэффициентами 


и правой частью ищется в виде интерполяционного полинома


\begin{equation}


\varphi_{nm}(s,\sigma )=\frac{4}{nm}\sum_{k=1}^n\ \sum_{i=1}^m


\alpha_{ki}\, \Delta_n(s-s_k)\Delta_m(\sigma -\sigma_i)\quad


(n,m\in \N),


\end{equation}


где $\Delta_r(t)$ --- обычное или модифицированное ядро Дирихле порядка


$\lbr r/2 \rbr$ при $r$ нечетном и соответственно при $r$ четном, где $r+1\in\N$.


Неизвестные коэффициенты $\alpha_{ki}=\varphi _{nm}(s_k,\sigma _i)$


определяются из следующей СЛАУ порядка $nm$:


\begin{multline}


a_0(s_j,\sigma _l)\alpha _{jl}+\frac{1}{n}\sum_{k=1}^n a_1(s_j,\sigma _l;


s_k)\beta _{j-k}(n)\alpha _{kl}+


\frac{1}{m}\sum_{i=1}^m a_2(s_j,\sigma _l;


\sigma _i)\beta _{l-i}(m)\alpha _{ji}+ \\


%


+\frac{1}{nm}\sum_{k=1}^n\;\sum_{i=1}^m a_{12}(s_j,\sigma _l;s_k,\sigma _i)


\beta _{j-k}(n)\beta _{l-i}(m)\alpha _{ki}=


f(s_j,\sigma _l);\ \ j=\overline{1,n},\ \ l=\overline{1,m},


\end{multline}


где


\begin{gather}


s_j=s_{jn}=\frac{2j\pi}{n},\  \ j=\overline{1,n};\quad


\sigma_l=\si_{lm}=\frac{2l\pi}{m},\ \ l=\overline{1,m}, \\


%


\beta _{j-k}(r)=\begin{cases}


\tg \dfrac{j-k}{2r}\pi , & \text{если } \ j-k \ \text{ четно}; \\[3mm]


\ctg \dfrac{k-j}{2r}\pi , &


\text{если }\ j-k \ \text{ нечетно},


\end{cases}


\end{gather}


при $r=2r_1+1$ ($r_1=0,1,\dotsc$), а при $r=2r_1$


($r_1=1,2,\dotsc$)


\begin{equation}


\beta _{j-k}(r)=\begin{cases}


0 , & \text{если } \ j-k \ \text{ четно}; \\


2\ctg \dfrac{k-j}{r}\pi , &


\text{если } \ j-k \ \text{ нечетно}.


\end{cases}


\end{equation}





Для вычислительной схемы (0.1), (1.27)--(1.31) справедлива





\begin{thm}


Если $f\in C[0,2\pi]^2$, то в условиях теоремы $1$ при любых


$n$ и $m\in \N$ СЛАУ $(1.28)$ имеет единственное решение


$\alpha _{jl}$ $(j=\overline{1,n}$, $l=\overline{1,m})$.


Приближенное решение $(1.27)$ удовлетворяет неравенству


\begin{equation}


\|\varphi_{nm}\|_{L_2[0,2\pi]^2}\leq \alpha ^{-1}


\|f\|_{C[0,2\pi]^2}\quad (n,m \in \N),


\end{equation}


является $L_2$-устойчивым


относительно малых возмущений правой части $f(s,\sigma)$ в


пространстве $C[0,2\pi ]^2$.


\end{thm}








{\bf Доказательство.} Обозначим через $X_{nm}$ множество всех двойных 


тригонометрических полиномов степени $(\lbr n/2\rbr,\lbr m/2 \rbr)$ 


вида (1.27), а через $\rL_{nm}:X\to X_{nm}\subset X$ --- оператор, 


ставящий в соответствие каждой функции $g\in C[0,2\pi]^2$ ее двумерный 


тригонометрический интерполяционный полином Лагранжа вида


$$


\rL_{nm}(g;s,\si) = \frac{4}{nm}\sum_{j=1}^{n}\ \sum_{l=1}^{m}


g(s_j,\si_l)\Delta_n(s-s_j)\Delta_m(\si-\si_l)


$$


по узлам (1.29). Тогда, как известно [4]--[8], СЛАУ (1.28) эквивалентна 


операторному уравнению


\begin{multline}


A_{nm}\vp_{nm}\equiv \rL_{nm}(a_0\vp_{nm})+\rL_{nm}S_1\rL_{n}^\xi(a_1\vp_{nm})+


\rL_{nm}S_2\rL_m^\eta(a_2\vp_{nm})+\\


+\rL_{nm}S_{12}\rL_{nm}^{\xi\eta}(a_{12}\vp_{nm})=\rL_{nm}f\quad


(\vp_{nm},\rL_{nm}f\in X_{nm}),


\end{multline}


где $\rL_n^\xi g$, $\rL_m^\eta g$ и $\rL_{nm}^{\xi\eta} g$ означают, 


что операторы $\rL_{n\infty}$, $\rL_{\infty m}$ и $\rL_{nm}$ применяются 


к непрерывной функции $g(s,\sigma;\xi,\eta)$ по переменным 


$\xi,\ \eta$ и $(\xi,\eta)$ соответственно.





Рассмотрим случай нечетных $n$ и $m \in \N$; другие случаи в силу 


результатов автора [4]--[10] рассматриваются аналогично.





Для любого полинома $\vp_{nm}\in X_{nm}$ с учетом свойств оператора $\rL_{nm}$


[4], [6], [13] имеем


\begin{multline}


(\rL_{nm}(a_0\vp_{nm}),\vp_{nm}) = \frac{1}{4\pi^2}\int_0^{2\pi}\int_0^{2\pi}


\rL_{nm}(a_0\vp_{nm};s,\si)\varphi_{nm}(s,\sigma)ds\, d\si = \\


%


=\frac{1}{nm}\sum_{j=1}^n\ \sum_{l=1}^m a_0(s_j,\si_l)


|\vp_{nm}(s_j,\si_l)|^2ds\, d\si \gea


\frac{\alpha}{nm}\sum_{j=1}^n\ \sum_{l=1}^m|\vp_{nm}(s_j,\si_l)|^2 = \\


%


=\frac{\alpha}{4\pi^2}\int_0^{2\pi}\int_0^{2\pi}|\vp_{nm}(s,\sigma)|^2ds\,d\si = 


\alpha \|\vp_{nm}\|^2\quad (\vp_{nm}\in X_{nm}).


\end{multline}


Для любого $\varphi_{nm}\in X_{nm}$ с учетом свойств использованных в 


работе квадратурных формул находим


\begin{gather*}


(\rL_{nm}S_1\rL_n^\xi (a_1\vp_{nm}),\vp_{nm}) = 


\frac{1}{4\pi^2}\int_0^{2\pi}\int_0^{2\pi}


\rL_{nm}[S_1\rL_n^\xi(a_1\vp_{nm});s,\si]\vp_{nm}(s,\si)ds\, d\si = \\


%


=\frac{1}{nm}\sum_{j=1}^n\ \sum_{l=1}^mS_1[\rL_n^\xi(a_1\vp_{nm});s_j,\si_l]


\vp_{nm}(s_j,\si_l) = \\


%


=\frac{1}{nm}\sum_{j=1}^n\ \sum_{l=1}^m \alpha_{jl}\frac{1}{n}


\sum_{k=1}^n a_1(s_j,\si_l;s_k)\beta_{j-k}(n)\vp_{nm}(s_k,\si_l) = \\


%


=\frac{1}{n^2m}\sum_{l=1}^m\bigg\{


\sum_{j=1}^n\alpha_{jl}\sum_{k=1}^na_1(s_j,\si_l;s_k)\beta_{j-k}(n)\alpha_{kl}


\bigg\} = \frac{1}{n^2m}\sum_{l=1}^m B_l,


\end{gather*}


где $\alpha_{jl} = \vp_{nm}(s_j,\si_l)$, $\alpha_{kl} = \vp_{nm}(s_k,\si_l)$.


В силу свойств функций $a_1(s,\si;\xi)$  и $\beta_{j-k}(n)$ величина


\begin{gather*}


B_l = \sum_{j=1}^n \alpha_{jl} \sum_{k=1}^na_1(s_j,\si_l;s_k)


\beta_{j-k}(n)\alpha_{kl} = \sum_{k=1}^n\alpha_{kl}


\sum_{j=1}^n a_1(s_k,\si_l;s_j)\beta_{k-j}(n)\alpha_{jl} = \\


%


=\sum_{j=1}^n \alpha_{jl}\sum_{k=1}^n a_1(s_j,\si_l;s_k)


[-\beta_{j-k}(n)]\alpha_{kl} = -B_l,\quad l=\ol{1,m}.


\end{gather*}


Cледовательно, $B_l=0$, $l=\ol{1,m}$, поэтому


\begin{equation}


(\rL_{nm}[S_1\rL_n^\xi(a_1\vp_{nm})],\vp_{nm}) = 0, \quad \vp_{nm}\in X_{nm}.


\end{equation}


Аналогично, с учетом свойств функций $a_2(s,\sigma;\eta)$ 


и $\beta_{l-i}(m)$ находим 


\begin{equation}


(\rL_{nm}[S_2\rL_m^\eta(a_2\vp_{nm})],\vp_{nm})=0, \quad \vp_{nm}\in X_{nm}.


\end{equation}


Кроме того, в силу свойств функций $a_{12}(s,\si;\xi,\eta)$


и $\beta_{j-k}(n)$, $\beta_{l-i}(m)$ для любого элемента\linebreak


$\vp_{nm}\in X_{nm}$


\begin{gather*}


(\rL_{nm}S_{12}\rL_{nm}^{\xi\eta}(a_{12}\vp_{nm}),\vp_{nm}) = 


\frac{1}{4\pi^2}\int_0^{2\pi}\int_0^{2\pi} \rL_{nm}[S_{12}\rL_{nm}(a_{12}


\vp_{nm});s,\si]\vp_{nm}(s,\si)ds\, d\si = \\


%


=\frac{1}{nm}\sum_{j=1}^n\ \sum_{l=1}^m \vp_{nm}(s_j,\si_l)S_{12}


[\rL_{nm}(a_{12}\vp_{nm});s_j,\si_l] = \\


%


=


\frac{1}{nm}\sum_{j=1}^n\ \sum_{l=1}^m \alpha_{jl}


\frac{1}{nm}\suml_{k=1}^n\sum_{i=1}^m


a_{12}(s_j,\si_l;s_k,\si_i)\beta_{j-k}(n)\beta_{l-i}(m)\alpha_{ki} = \\


%


=\frac{1}{nm}\sum_{k=1}^n\ \sum_{i=1}^m \alpha_{ki}\frac{1}{nm}


\suml_{j=1}^n\sum_{l=1}^m


a_{12}(s_k,\si_i;s_j,\si_l)\beta_{k-j}(n)\beta_{i-l}(m)\alpha_{jl} = \\


%


%%=\frac{1}{nm}\suml_{j=1}^n\suml_{l=1}^m \alpha_{jl} \frac{1}{nm}


%%\suml_{k=1}^n\suml_{i=1}^n[-a_{12}(s_j,\si_l;s_k,\si_i)]


%%\beta_{j-k}(n)\beta_{l-i}(m)\alpha_{ki} =


%


=-\frac{1}{nm}\sum_{j=1}^n\ \sum_{l=1}^m \alpha_{jl}\frac{1}{nm}


\suml_{k=1}^n\sum_{i=1}^n


a_{12}(s_j,\si_l;s_k,\si_i)\beta_{j-k}(n)\beta_{l-i}(m)\alpha_{ki}.


\end{gather*}


Поэтому для любого полинома $\vp_{nm}\in X_{nm}$ имеем 


\begin{equation}


(\rL_{nm}S_{12}\rL_{nm}^{\xi\eta}(a_{12}\vp_{nm}),\vp_{nm})=0.


\end{equation}





Из соотношений (1.33)--(1.37) следует


$$


(A_{nm}\vp_{nm},\vp_{nm})\gea \alpha\|\vp_{nm}\|^2, \quad \vp_{nm}\in X_{nm}.


$$


Отсюда для любого полинома $\vp_{nm}\in X_{nm}$ имеем неравенство


$$


\|A_{nm}\vp_{nm}\|\gea \alpha\|\vp_{nm}\|,\quad \vp_{nm}\in X_{nm},


$$


обеспечивающее существование двусторонне обратного оператора 


$A_{nm}^{-1}:X_{nm}\to X_{nm}$ и справедливость неравенств 


(напр., [13], гл.\,I, \S\,2)


\begin{equation}


\|A_{nm}^{-1}\|\lea \alpha^{-1}<\infty\quad (n,m\in \N).


\end{equation}


Поэтому операторное уравнение (1.33) и эквивалентная ему СЛАУ (1.28) 


однозначно разрешимы при любых $n,m \in \N$ и любых правых частях, 


причем $\vp_{nm} = A_{nm}^{-1}\rL_{nm}f$. Отсюда и из (1.38) находим неравенства


\begin{gather*}


\|\vp_{nm}\|_{L_2}\lea \|A_{nm}^{-1}\|\|\rL_{nm}f\|\lea \frac{1}{\alpha}\bigg\{


\frac{1}{4\pi^2}\int_0^{2\pi}\int_0^{2\pi}|\rL_{nm}(f;s,\si)|^2ds\, d\si


\bigg\}^{1/2}=\\


%


=\frac{1}{\alpha}\bigg\{\frac{1}{nm}\sum_{j=1}^n\ \sum_{l=1}^m|f(s_j,\si_l)|^2


\bigg\}^{1/2}\lea \frac{1}{\alpha}\|f\|_{C}.


\end{gather*}





Oстальные утверждения теоремы 5 легко выводятся из неравенства (1.32).





Далее, обозначим через $\overline{X}=\{ \overline{\varphi }\}$


пространство всех векторов вида


$\overline{\varphi }=\{\alpha _{jl}\}_{j=\ol{1,n}}^{l=\ol{1,m}}$,


где $\alpha _{jl}\in \R$, $j=\overline{1,n}$, $l=\overline{1,m}$,


с нормой


$$


\|\overline{\varphi}\|_{\overline{X}}=\bigg\{\frac{1}{nm}


\sum_{j=1}^n\sum_{l=1}^m|\alpha_{jl}|^2\bigg\}^{1/2}\quad


(\ol{\varphi}\in \overline{X})


$$


и со скалярным произведением


$$


(\ol{\vp},\ol{\psi}) = \frac{1}{nm}\sum_{j=1}^n\ \sum_{l=1}^m


\alpha_{jl}\beta_{jl}\quad (\ol{\vp},\ol{\psi}\in \ol{X}),


$$


где $\ol{\vp}=\{\alpha_{jl}\}_{j=\ol{1,n}}^{l=\ol{1,m}}$ и 


$\ol{\psi} = \{\beta_{jl}\}_{j=\ol{1,n}}^{l=\ol{1,m}}$.


Тогда СЛАУ (1.28) можно записать в виде эквивалентного ей линейного уравнения


\begin{equation}


\overline{A}\overline{\varphi }=\overline{f}\quad 


(\overline{\varphi }, \overline{f}\in \overline{X}),


\end{equation}


где $\overline{f}=\{ f(s_j,\sigma _l)\}_{j=\ol{1,n}}^{l=\ol{1,m}}$, 


а $\overline{A}:\overline{X}\to \overline{X}$ --- матрица, определяемая 


левой частью СЛАУ (1.28).





Из вышеприведенных теорем 1, 2, 5 и из результатов


([13], гл.\,II, \S\,7) выводится





\begin{thm}


В условиях теоремы $5$ решение 


$\overline{\varphi }=\{\alpha _{jl}\}_{j=\ol{1,n}}^{l=\ol{1,m}}\in


\overline{X}$ СЛАУ $(1.28)$ можно найти с помощью универсального


итерационного метода


\begin{equation}


\overline{\varphi }^k=\overline{\varphi }^{k-1}+\alpha \, \|\overline{A}\|^{-2}


(\overline{f}-\overline{A}\overline{\varphi }^{k-1}),\quad k=1,2,\dotsc,


\end{equation}


при любом начальном приближении 


$\overline{\varphi}^0=\{\alpha_{jl}^0\}_{j=\ol{1,n}}^{l=\ol{1,m}}


\in \overline{X}$. Погрешность $k$-го приближения


$\overline{\varphi }^k=\{\alpha _{jl}^k\}_{j=\ol{1,n}}^{l=\ol{1,m}}


\in \overline{X}$ может быть


оценена в пространстве $\overline{X}$ неравенствами


\begin{equation}


\|\overline{\varphi }-\overline{\varphi }^k\|_{\overline{X}}=


\bigg\{\frac{1}{nm}


\sum_{j=1}^n\ \sum_{l=1}^m|\alpha _{jl}-\alpha _{jl}^k|^2\bigg\}^{1/2}\lea


\overline{q}^k(1-\overline{q})^{-1} \|\overline{\varphi }^1-


\overline{\varphi }^0\|_{\overline{X}},


\quad k=1,2,\dotsc,


\end{equation}


где $\overline{q}=(1-\alpha ^2\|\overline{A}\|^{-2})^{1/2}<1;$


если же начальное  приближение определяется по формуле 


$\overline{\varphi }^0=\alpha \|\overline{A}\|^{-2}\overline{f}


\in \overline{X}$, то и неравенствами


\begin{equation}


\|\overline{\varphi }-\overline{\varphi }^k\|_{\overline{X}}


\leq \overline{q}^{k+1}(1-\overline{q})^{-1} \alpha


\|\overline{A}\|^{-2}\|\overline{f}\|_{\overline{X}},


\; k=1,2,\dotsc


\end{equation}


\end{thm}





Рассмотрим еще одну вычислительную схему полиномиального 


квадратурно-кубатурного метода решения бисингулярного интегрального 


уравнения (0.1). Приближенное решение снова будем искать 


в виде тригонометрического интерполяционного полинома (1.27) по узлам (1.29), 


коэффициенты $\alpha_{kj} = \alpha_{kj}(n,m)$ которого определим из СЛАУ 


\begin{gather}


a_0(s_j,\si_l)\alpha_{jl}+\frac{1}{n}


\sum_{\overset{k=1}{k\ne j}}^n a_1(s_j,\si_l;s_k)\ctg


\frac{(k-j)\pi}{n}\alpha_{kl}+


\frac{1}{m}\sum_{\overset{i=1}{i\ne l}}^m a_2(s_j,\si_l;\si_i)\ctg


\frac{(i-l)\pi}{m}\alpha_{ji}+ \notag \\


%


+\frac{1}{nm}\sum_{\overset{k=1}{k\ne j}}^{n}


\ \sum_{\overset{i=1}{i\ne l}}^m a_{12}(s_j,\si_l;s_k,\si_i)\ctg


\frac{(k-j)\pi}{n}\ctg \frac{(i-l)\pi}{m}\alpha_{ki}= \notag \\


=f(s_j,\si_l);\ \ j=\ol{1,n},\ \ l=\ol{1,m}\ \ (n,m,\in \N).


\end{gather}





Для вычислительной схемы (0.1), (1.27), (1.29), (1.43) справедливы 


результаты, аналогичные теоремам 5 и 6. Например, имеет место





\begin{thm}


Пусть $a_0(s,\si)$, $a_1(s,\si,\xi)$, $a_2(s,\si;\eta)$, 


$a_{12}(s,\si;\xi,\eta)$ и $f(s,\si)$ --- непрерывные $2\pi$-периодические 


функции по каждой из переменных.  


Тогда в условиях теоремы $1$ СЛАУ $(1.43)$ однозначно разрешима при любых 


$n$ и $m\in \N$, а приближенные решения $(1.27)$ удовлетворяют неравенствам


\begin{equation}


\|\vp_{nm}\|_{L_2}\lea \frac{1}{\alpha}\bigg\{\frac{1}{nm}


\sum_{j=1}^n\ \sum_{l=1}^m|f(s_j,\si_l)|^2\bigg\}^{1/2}\lea


\frac{1}{\alpha}\|f\|_C\quad (n,m\in\N).


\end{equation}


\end{thm}





\begin{pf}


Представим СЛАУ (1.43) в операторном виде (1.39), где 


$\ol{A}:\ol{X}\to \ol{X}$ --- матричный оператор, определяемый левой 


частью СЛАУ (1.43). Тогда для любого вектора 


$\ol{\vp}=\{\alpha_{jl}\}_{j=\ol{1,n}}^{l=\ol{1,m}}$ находим


\begin{multline}


(\ol{A}\ol{\vp},\ol{\vp}) = \frac{1}{nm}\sum_{j=1}^n\ \sum_{l=1}^m


\alpha_{jl}\bigg\{a_0(s_j,\si_l)\alpha_{jl}+


\frac{1}{n}\sum_{\overset{k=1}{k\ne j}}^n \alpha_{kl}a_1(s_j,\si_l;s_k)\ctg


\frac{s_k-s_j}{2}+\\ 


%


+\frac{1}{m}\sum_{\overset{i=1}{i\ne l}}^m\alpha_{ji}


a_2 (s_j,\si_l;\si_i)\ctg\frac{s_i-s_l}{2}+ \\


%


+


\frac{1}{nm}\sum_{\overset{k=1}{k\ne j}}^n


\ \sum_{\overset{i=1}{i\ne l}}^m


\alpha_{ki}a_{12}(s_j,\si_l;s_k,\si_i)\ctg\frac{s_k-s_j}{2}


\ctg\frac{s_i-s_l}{2}\bigg\}\gea \\


%


\gea\alpha\frac{1}{nm}\sum_{j=1}^n\ \sum_{l=1}^m|\alpha_{jl}|^2+M_1+M_2+M_{12} = 


\alpha\|\ol{\vp}\|^2_X+M_1+M_2+M_{12},


\end{multline}


где смысл обозначений $M_i$ очевиден.


С учетом свойств функций $a_1(s,\si;\xi)$ и $\ctg\frac{\xi-s}{2}$ находим


\begin{multline}


M_1=\frac{1}{nm}\sum_{j=1}^n\ \sum_{l=1}^m \alpha_{jl}\frac{1}{n}


\suml_{\overset{k=1}{k\ne j}}^n\alpha_{kl}a_1(s_j,\si_l;s_k)


\ctg\frac{s_k-s_j}{2} = \\


%


=\frac{1}{n^2m}\suml_{l=1}^m\bigg[


\sum_{j=1}^n\alpha_{jl}\sum_{\overset{k=1}{k\ne j}}^n \alpha_{kl}


a_1(s_j,\si_l;s_k)\ctg\frac{s_k-s_j}{2}\bigg]= \\


%


=\frac{1}{n^2m}\sum_{l=1}^m\bigg[\sum_{k=1}^n\alpha_{kl}


\sum_{\overset{j=1}{j\ne k}}^n\alpha_{jl}a_1(s_k,\si_l;s_j)


\ctg\frac{s_j-s_k}{2}\bigg]= \\


%


=\frac{1}{n^2m}\sum_{l=1}^m\bigg[


\sum_{j=1}^n\alpha_{jl}\sum_{\overset{k=1}{k\ne j}}^n


\alpha_{kl}a_1(s_j,\si_l;s_k)


\bigg(-\ctg \frac{s_k-s_j}{2}\bigg)\bigg]=-M_1=0.


\end{multline}





Аналогично для любого вектора $\ol{\vp}=\{\alpha_{jl}\}\in \ol{X}$ находим


\begin{equation}


M_2 = \frac{1}{nm}\sum_{j=1}^n\ \sum_{l=1}^m \alpha_{jl}\frac{1}{m}


\sum_{\overset{i=1}{i\ne l}}^m \alpha_{ji}a_2(s_j,\si_l;\si_i)


\ctg \frac{s_i-s_l}{2}=0.


\end{equation}


С учетом свойств функций $h_{12}(s,\si;\xi,\eta)$  и $\ctg\frac{\xi-s}{2}$, 


$\ctg\frac{\eta-\si}{2}$ для любого $\ol{\vp}\in \ol{X}$ получаем


\begin{equation}


M_{12} = \frac{1}{nm}\sum_{j=1}^{n}\ \sum_{l=1}^m\alpha_{jl}\frac{1}{nm}


\sum_{\overset{k=1}{k\ne j}}^n\ \sum_{\overset{i=1}{i\ne l}}^m


\alpha_{ki}h_{12}(s_j,\si_l;s_k,\si_i)\ctg\frac{s_k-s_j}{2}


\ctg\frac{s_i-s_l}{2}=0.


\end{equation}





Из соотношений (1.45)--(1.48) находим неравенство


\begin{equation}


(\ol{A}\ol{\vp},\ol{\vp})\gea \alpha\|\ol{\vp}\|^2_{\ol{X}}\quad


(\ol{\vp}\in \ol{X}).


\end{equation}





Из (1.49) следует, что оператор $\ol{A}:\ol{X}\to \ol{X}$ имеет двусторонний 


обратный $\ol{A}{} ^{-1}$ и 


$$


\|\ol{A}{}^{-1}\|\lea \alpha^{-1}<\infty,\quad 


\ol{A}{}^{-1}:\ol{X}\to \ol{X}.


$$


Поэтому СЛАУ (1.43) имеет единственное решение 


$\ol{\vp} = \{\alpha_{jl}\}_{j=\ol{1,m}}^{l=\ol{1,m}}\in \ol{X}$ и 


для него справедливы соотношения


\begin{equation}


\|\ol{\vp}\|_{\ol{X}} = \bigg\{


\frac{1}{nm}\sum_{j=1}^n\;\sum_{l=1}^m|\alpha_{jl}|^2\bigg\}^{1/2}\lea


\frac{1}{\alpha}\bigg\{


\frac{1}{nm}\sum_{j=1}^n\;\sum_{l=1}^m|f(s_j,\si_l)|^2


\bigg\}^{1/2}


\equiv \frac{1}{\alpha}\|\ol{f}\|_{\ol{X}}\ \ (n,m\in \N).


\end{equation}


Из соотношений (1.27), (1.29) и (1.50) с помощью соответствующих 


результатов [4], [6], [10] находим оценки (1.44).


\end{pf}





1.5.2. {\it Сплайновые квадратурно-кубатурные методы.} Теперь приближенное 


решение с.\,и.\,у. (0.1) будем искать в виде двумерного сплайна нулевой степени


\begin{equation}


\vp_{nm}(s,\si) = \sum_{k=1}^n\ \sum_{i=1}^m


\alpha_{ki}\psi_{kn}(s)\psi_{im}(\si) \quad (n,m\in \N),


\end{equation}


где $\psi_{kn}(s)$ и $\psi_{im}(\si)$, как и выше, суть 


$2\pi$-периодические характеристические функции интервалов 


соответственно $[s_{k-1},s_k)$ и $[\si_{i-1},\si_i)$, а узлы определены 


в (1.23). Неизвестные постоянные $\alpha_{ki} = \alpha_{ki}(n,m)$ будем 


определять из СЛАУ


\begin{multline}


a_0(\ol{s}_j,\ol{\si}_l)\alpha_{jl}+\sum_{k=1}^n \alpha_{kl}


a_1(\ol{s}_j,\ol{\si}_l;s_k)


\gamma_{k-j}{(n)}+\sum_{i=1}^m\alpha_{ji}a_2


(\ol{s}_j,\ol{\si}_l;\ol{\si}_i)\gamma_{i-l}{(m)}+ \\


%


+ \sum_{k=1}^n\ \sum_{l=1}^m\alpha_{ki}a_{12}(\ol{s}_j,


\ol{\si}_l;\ol{s}_k,\ol{\si}_i)\gamma_{k-j}{(n)}\gamma_{i-l}{(m)}=


f(\ol{s}_j,\ol{\si}_l);\quad j=\ol{1,n},\ \ l=\ol{1,m},


\end{multline}


где 


\begin{gather}


\ol{s}_j = \frac{(2j-1)\pi}{n},\quad \ol{\sigma}_l = \frac{(2l-1)\pi}{m}, \\


%


\gamma_{k-j}(n) = \frac{1}{2\pi}\ln \left|


\frac{\sin s_{k-j+1/2}}{\sin s_{k-j-1/2}}\right|,


\quad \gamma_{i-l}(m) = \frac{1}{2\pi}\ln


\left|\frac{\sin s_{i-l+1/2}}{\sin s_{i-l-1/2}}\right|.


\end{gather}





Для вычислительной схемы (0.1), (1.23), (1.51)--(1.54) справедлива 





\begin{thm}


В условиях теоремы $7$ СЛАУ $(1.52)$--$(1.54)$ однозначно разрешима при 


любых $n$ и $m\in \N$, а для приближенного решения $(1.51)$ справедливы 


соотношения


$$


\|\vp_{nm}\|_{L_2} = \left\{


\frac{1}{nm}\sum_{k=1}^n\ \sum_{i=1}^m|\alpha_{ki}|^2


\right\}^{1/2}\lea \frac{1}{\alpha}\left\{


\frac{1}{nm}\sum_{j=1}^n\ \sum_{l=1}^m |f(\ol{s}_j,\ol{\si}_l)|^2


\right\}^{1/2} \quad (n,m\in \N).


$$


\end{thm}





{\bf Доказательство} можно вести как по схеме доказательства теоремы 5, 


так и по схеме доказательства теоремы 7; при этом существенно используются 


свойства функций $a_i$ ($i=0,1,2,12$), $\gamma_{k-j}(n)$,


$\gamma_{i-l}(m)$ и ядер Гильберта


$\ctg\frac{\xi-s}{2}$, $\ctg\frac{\eta-\si}{2}$. Ввиду громоздкости 


на деталях останавливаться не будем, тем более они почти аналогичны 


одномерному случаю [14]. 





Теперь приближенное решение с.\,и.\,у. (0.1) будем искать в виде 


двумерного сплайна первой степени вида


\begin{equation}


\widetilde{\vp}_{nm}(s,\si) = \sum_{k=1}^n\ \sum_{i=1}^m


\alpha_{ki}\psi_{kn}(s) \psi_{im}(\si)\quad (n,m\in \N),


\end{equation}


где $\psi_{kn}(s)$ и $\psi_{im}(\si)$ суть одномерные $2\pi$-периодические 


фундаментальные сплайны первой степени по системам узлов  


$s_k = 2k\pi/n$ ($k=\ol{0,n}$) и $\si_i = 2i\pi/m$ ($i=\ol{0,m}$) соответственно,


$$


\psi_{kn}(s)=


\begin{cases}


0 & \text{при } s \lea s_{k-1};\\


\dfrac{s-s_{k-1}}{s_k - s_{k-1}} & \text{при } s_{k-1}\lea s \lea s_k;\\


\dfrac{s_{k+1}-s}{s_{k+1} - s_{k}} & \text{при } s_{k}\lea s \lea s_{k+1};\\


0 & \text{при } s \gea s_{k+1},


\end{cases}


\quad


\psi_{im}(\si)=


\begin{cases}


0 & \text{при } \si \lea \si_{i-1}; \\ 


\dfrac{\si-\si_{i-1}}{\si_i - \si_{i-1}} & \text{при } \si_{i-1}\lea


\si \lea \si_i; \\


\dfrac{\si_{i+1}-\si}{\si_{i+1} - \si_{i}} & \text{при } \si_{i}\lea


\si \lea \si_{i+1}; \\


0 & \text{при } \si \gea \si_{i+1}.


\end{cases}


$$


Неизвестные коэффициенты $\alpha_{ki}=\alpha_{ki}(n,m)$ будем определять из СЛАУ


\begin{gather}


a_0(s_j,\si_l)\alpha_{jl}+\sum_{k=1}^n \alpha_{kl}a_1(s_j,\si_l;s_k)


\delta_{k-j}(n)+


\sum_{i=1}^m\alpha_{ji}a_2(s_j,\si_l;\si_i)\delta_{i-l}(m)+\notag \\


%


+\sum_{k=1}^n\; \sum_{i=1}^m \alpha_{ki} a_{12}(s_j,\si_l;s_k,\si_i)


\delta_{k-j}(n)\delta_{i-l}(m) = f(s_j,\si_l),\quad j=\ol{1,n},\; l=\ol{1,m},


\end{gather}


где узлы приведены в (1.23), а


\begin{align}


\delta_{k-j}(n) &= \frac{1}{n}\int_0^1\left[\ctg


\frac{\pi}{n}(k-j+t)+\ctg \frac{\pi}{n}(k-j-t)\right](1-t)dt, \\


%


\delta_{i-l}(m) &= \frac{1}{m}\intl_0^1\left[


\ctg \frac{\pi}{m}(i-l+s)+\ctg\frac{\pi}{m}(i-l-s)\right](1-s)ds.


\end{align}





Для вычислительной схемы (0.1), (1.23), (1.55)--(1.58) справедлива 





\begin{thm}


В условиях теоремы $7$ СЛАУ $(1.56)$--$(1.58)$ однозначно разрешима при 


любых $n$ и $m \in \N$, а для приближенного решения $(1.55)$ справедливы 


неравенства


$$


\|\widetilde{\vp}_{nm}\|\lea \left\{


\frac{1}{nm}\sum_{k=1}^n\;\sum_{i=1}^m|\alpha_{ki}|^2\right\}^{1/2}\lea 


\frac{1}{\alpha}\left\{


\frac{1}{nm}\sum_{j=1}^n\;\sum_{l=1}^m|f(s_j,\sigma_l)|^2\right\}^{1/2}


\quad (n,m\in \N).


$$


\end{thm}





{\bf Доказательство} проводится по схеме доказательства теоремы 7


с использовнием соответствующего результата для одномерного 


случая ([10], с.\,152--153); ввиду громоздкости на подробных выкладках 


останавливаться не будем.








\section{Некоторые замечания}








\begin{note}


Теоремы 1--4 остаются в силе при возмущении оператора $A:L_2\to L_2$


соответствующим оператором $T$: а)~вполне непрерывным в $L_2=L_2[0,2\pi]^2$;


б)~малым по норме оператором $T:L_2 \to L_2$ таким, что


$\|T\|<\alpha$, где постоянная $\alpha$ определена в теореме~1.


\end{note}








\begin{note}


Теоремы 5--9 остаются в силе при возмущении оператора $A:L_2\to L_2$


соответствующим оператором $T$: а)~вполне непрерывным из пространства


$L_2$ в пространство $C=C[0,2\pi]^2$;


б)~малым по норме в том смысле, что


$\|T\|_{{}_{\scriptstyle L_2\rightarrow C}}<\alpha$.


\end{note}





\begin{note}


На основе приведенных выше теорем 1--9 и результатов [13]--[16]


могут быть построены быстросходящиеся в $L_2[0,2\pi ]^2$


схемы кубатурно-итерационных методов решения


бисингулярного интегрального уравнения (0.1).


\end{note}








\begin{note}


В формулах (1.6)--(1.8), (1.26) и соответственно (1.40)--(1.42) нормы


$\|A\|$ и $\| \overline{A}\|$ могут быть заменены постоянными


$M$ и $\overline{M}$ соответственно, где $\|A\|\lea


M< \infty$, $A:X\to X,$ и $\|\overline{A}\|\lea \overline{M}<\infty$,


$\ol{A}:\ol{X}\to \ol{X}$, причем $\ol{M}\lea M,$ а постоянная $M$,


как указано выше, определяется через функции $a_i$ ($i=0,1,2,12$).


\end{note}





\begin{thebibliography}{99}





\bibitem{1}


Гахов Ф.Д. {\em Краевые задачи}. -- М.: Наука, 1977. -- 638~с.


\bibitem{2}


Чибрикова Л.И. {\em Основные граничные задачи для аналитических функций}.


-- Казань: Изд-во Казанск. ун-та, 1977. -- 302~с.


\bibitem{3}


Michlin S.G., Pr\"{o}\ss{}dorf S. {\em Singul\"{a}re Integraloperatoren}.


-- Berlin: Akademic-Verlag, 1980. -- 514~S.


\bibitem{4}


Габдулхаев Б.Г. {\em Кубатурные формулы для многомерных сингулярных интегралов},


I {\it и} II // Тр. Матем. ин-та АН Болгарии. -- 1970. -- Т.\,11. --


С.\,181--196; Изв. вузов. Математика. -- 1975. -- \No\,4. -- С.\,3--13.


\bibitem{5}


Габдулхаев Б.Г. {\em Прямые методы решения некоторых операторных уравнений},


I--IV // Изв. вузов. Математика. -- 1971. -- \No\,11. -- С.\,33--44;


1971. -- \No\,12. -- С.\,28--38; 1972. -- \No\,4. -- С.\,32--43;


1974. -- \No\,3. -- С.\,18--31.


\bibitem{6}


Габдулхаев Б.Г. {\em Приближенное решение многомерных сингулярных 


интегральных уравнений}, I {\it и} II // Изв. вузов. Математика. -- 


1975. -- \No\,7. -- С.\,30--41; 1976. -- \No\,1. -- С.\,30--41.


\bibitem{7}


Габдулхаев Б.Г. {\em Конечномерные аппроксимации сингулярных интегралов 


и прямые методы решения особых интегральных и интегродифференциальных 


уравнений} // Итоги науки и техн. Матем. анализ. -- М.: Изд-во АН СССР, 


1980. -- Вып. 18. -- С.\,251--307.


\bibitem{8}


Габдулхаев Б.Г. {\em Многомерные сингулярные интегральные уравнения с 


положительными операторами} // Дифференц. уравнения. -- 1993. -- Т.29.


-- \No\,9. -- С.\,1504--1516.


\bibitem{9}


Габдулхаев Б.Г. {\em Конечномерные приближения решений бисингулярного 


интегрального уравнения} // Теория функций и ее прилож. Тез. докл. 


Школы-конф., 15--22 июня 1995~г. -- Казань: Казанский 


фонд ``Математика'', 1995. -- С.\,15-17.


\bibitem{10}


Габдулхаев Б.Г. {\em Численный анализ сингулярных интегральных уравнений. 


Избранные главы}. -- Казань: Изд-во Казанск. ун-та, 1995. -- 230~с.


\bibitem{11}


Канторович Л.В.,  Акилов Г.П. {\em Функциональный анализ в нормированных 


пространствах}. -- М.: Физматгиз, 1959. -- 684~с.


\bibitem{12}


Гаевский Х., Грeгер К., Захариас К. {\em Нелинейные операторные уравнения 


и операторные дифференциальные уравнения}. -- М.: Мир, 1978. -- 336~с.


\bibitem{13}


Габдулхаев Б.Г. {\em Оптимальные аппроксимации решений линейных задач}.


-- Казань: Изд-во Казанск. ун-та, 1980. -- 232~с.


\bibitem{14}


Габдулхаев Б.Г. {\em Методы решения сингулярных интегральных уравнений с 


положительными операторами} // Дифференц. уравнения. -- 1997. -- Т.\,33.


-- \No\,3. -- С.\,400--409.


\bibitem{15}


Лучка А.Ю. {\em Проекционно-итеративные методы}. -- Киев: Наук. думка,


1993. -- 288~с.


\bibitem{16}


Габдулхаев Б.Г. {\em Решение операторных уравнений методом уточняющих итераций}


// Изв. вузов. Математика. -- 1974. -- \No\,5. -- С.\,66--80.





\end{thebibliography}





\vskip 2cm





\post{\it Казанский государственный}{\it Поступила}


\post{\it университет}{20.05.2003}





\label{end}


\end{document}





